\documentclass[
	a4paper,
	% showframes,
	% vline=2.2em,
	% maincolor=cvgreen,
	% sidecolor=gray!50,
	% sectioncolor=red,
	% subsectioncolor=orange,
	% itemtextcolor=black!80,
	% sidebarwidth=0.4\paperwidth,
	% topbottommargin=0.03\paperheight,
	% leftrightmargin=20pt,
	% profilepicsize=4.5cm,
	% profilepicborderwidth=3.5pt,
	% profilepicstyle=profilecircle,
	% profilepiczoom=1.0,–
	% profilepicxshift=0mm,
	% profilepicyshift=0mm,
	% profilepicrounding=1.0cm,
]{cv_class}

% improve word spacing and hyphenation
\usepackage{microtype}
\usepackage{ragged2e}

% uncomment in case you don't want any hyphenation
% \usepackage[none]{hyphenat}

% take care of proper font encoding
\ifxetexorluatex
	\usepackage{fontspec}
	\defaultfontfeatures{Ligatures=TeX}
%	\newfontfamily\headingfont[Path = fonts/]{segoeuib.ttf} % local font
\else
	\usepackage[utf8]{inputenc}
	\usepackage[T1]{fontenc}
%	\usepackage[sfdefault]{noto} % use noto google font
\fi

% enable mathematical syntax for some symbols like \varnothing
\usepackage{amssymb}

% bubble diagram configuration
\usepackage{smartdiagram}
\smartdiagramset{
	% default font size is \large, so adjust to harmonize with sidebar layout
	bubble center node font = \footnotesize,
	bubble node font = \footnotesize,
	% default: 4cm/2.5cm; make minimum diameter relative to sidebar size
	bubble center node size = 0.4\sidebartextwidth,
	bubble node size = 0.25\sidebartextwidth,
	distance center/other bubbles = 1.5em,
	% set center bubble color
	bubble center node color = maincolor!70,
	% define the list of colors usable in the diagram
	set color list = {maincolor!10, maincolor!40,
	maincolor!20, maincolor!60, maincolor!35},
	% sets the opacity at which the bubbles are shown
	bubble fill opacity = 0.8,
}


%-------------------------------------------------------------------------------
%                            PERSONAL INFORMATION
%-------------------------------------------------------------------------------
%% mandatory information
% your name
\cvname{Даниил \\ Цыганков}
% job title/career
\cvjobtitle{Врач-стоматолог-хирург}

%% optional information
% profile picture
\cvprofilepic{pics/profile_4.jpg}

% NOTE: ordering in sidebar will mimic the following order
% date of birth
\cvbirthday{Мужчина, 25 лет}
% short address/location, use \newline if more than 1 line is required
% phone number
\cvphone{+7 (987) 690-60-86}
% personal website
% email address
\cvmail{dani1l.tsygankov@yandex.ru}
% pgp key
% any other custom entry
\cvaddress{Россия,  Москва,  м.Лухмановская}

%-------------------------------------------------------------------------------
%                              SIDEBAR 1st PAGE
%-------------------------------------------------------------------------------
% add more profile sections to sidebar on first page
\addtofrontsidebar{
	% include gosquare national flags from https://github.com/gosquared/flags;
	% naming according to ISO 3166-1 alpha-2 country codes
	\graphicspath{{pics/flags/}}

	% social network accounts incl. proper hyperlinks
	\profilesection{Контакты}
		\begin{icontable}{2.5em}{1em}
			\social{\faTelegram}
				{https://t.me/79876906086}
				{Telegram }
			\social{\faWhatsapp}
				{https://wa.me/79876906086}
				{WhatsApp}
		\end{icontable}

	\profilesection{Знание языков}
		\pointskill{\flag{RU.png}}{Русский (родной)}{5}
	\pointskill{\flag{GB.png}}{Английский (средний)}{3}

	\profilesection{Навыки}
		\skill{\faCompress}{Хирургическое лечение}
		\skill{\faCompress}{Диагностика}
		\skill{\faCompress}{Составление плана лечения}
		\skill{\faCompress}{Шинирование}
		\skill{\faCompress}{Работа в 4-ре руки}
		\skill{\faCompress}{Мед. документация}
		

	\profilesection{Личные качества}
		\skill{\faCompress}{Ответственность}
		\skill{\faCompress}{Эмоциональная устойчивость}
		\skill{\faCompress}{Сдержанность}
		\skill{\faCompress}{Трудолюбие}
}



%-------------------------------------------------------------------------------
%                         TABLE ENTRIES RIGHT COLUMN
%-------------------------------------------------------------------------------
\begin{document}

\makefrontsidebar

\cvsection{Обо мне}
Уважаемый читатель, меня зовут Даниил Александрович,  и сейчас я ищу работу врача стоматолога - хирурга. У меня высшее образование по специальности "Стоматология хирургическая". \\ Стаж работы в медицине 7 лет, работал в нескольких крупных стоматологических клиниках.  Заинтересован в постоянном обучении и личностном развитиии.  Ищу работу, чтобы получить новые навыки и попробовать свои силы в решении \\ сложных задач.

\cvsection{Опыт работы}
\begin{cvtable}[2]
				\cvitem{2021-- 2022}{Стоматологическая клиника "Darus dent" }{}{\textbf{Должность:} Врач-стоматолог-хирург.
\\\textbf{Обязанности:} Удаление зубов, шинирование, вскрытие поднадкостничного абсцесса, первичная консультация и составление плана лечения.}

		\cvitem{2019-- 2021}{ООО "Центр стоматологии" }{}{\textbf{Должность:} Врач-стоматолог-гигиенист\\\textbf{Обязанности:} Проведение первичного стоматологического осмотра, составление предварительного плана лечения, проведение профессиональной гигиены полости рта, подбор индивидуальных средств гигиены, направление к смежным специалистам. 
}

	\cvitem{2018 -- 2019}{ООО "Центр стоматологии" }{}{\textbf{Должность:} Ассистент стоматолога\\\textbf{Обязанности:} Работа в 4 руки, ведение медицинской документации, стерилизация инструментов, ассистирование на имплантациях и челюстно-лицевых операциях. }

	\cvitem{2015 -- 2018}{ГБУЗ РМ Мордовская республиканская центральная клиническая больница. }{}{\textbf{Должность:} Санитар приемного отделения №1.\\\textbf{Обязанности:} Транспортировка пациентов, ПХО ран, внутривенные/внутримышечные инъекции.}



	
\end{cvtable}


\cvsection{Образование}

\begin{cvtable}[2]


	\cvitem{2020 -- 2022}{ФГБОУ "НИМГУ им. Н.П. Огарева"}{}
		{Диплом об образовании и присвоении квалификации\\ \textbf{Врач-стоматолог-хирург}}
		
	\cvitem{2015 -- 2020}{ФГБОУ "НИМГУ им. Н.П. Огарева"}{}
		{Диплом об образовании и о квалификации\\ \textbf{Врач - стоматолог}}
		
\end{cvtable}


\cvsection{Публикации и награды}
\begin{cvtable}
	\cvitem{2020}{\textbf{Эволюция подходов к выбору метода хирургического лечения паховых грыж.} Филяева А.С., Цыганков Д.А. и др. \\XV Международная / XXVIV Всероссийская Пироговская научная медицинская конференция студентов и молодых ученых.}{}{}
	\cvitem{2019}{\textbf{Диплом Лауреата I степени} в номинации: Научные статьи по медицинским наукам в Международном Конкурсе Научных Работ PTScience.}{}{}
\end{cvtable}



\end{document}
